\documentclass[Total.tex]{subfiles}
\begin{document}
	\chapter{Rings}
		\section{Definitions}
		\section{UFD,PID,Euclidean Ring}
		\section{Localization}
		\section{title}
			\subsection{Ring Maps of Finite Type and of Finite Representation}
			\begin{Def}
				Let $R\rightarrow S$ be a ring map:
				\begin{itemize}
					\item $R\rightarrow S$ is \textbf{of finite type}, or that $S$ is a \textbf{finite type $R$-algebra} if there exist an $n\in \NN$ and an \underline{surjection} of $R$-algebras
					\begin{gather*}
						R[x_1,\dots,x_n]\rightarrow S
					\end{gather*}
					\item $R\rightarrow S$ is \textbf{of finite presentation} if there exist integers $n,m\in \NN$ and polynomials
					\begin{gather*}
						f_1,\dots,f_m\in R[x_1,\dots,x_n]
					\end{gather*}
					and an isomorphism of $R$-algebras
					\begin{gather*}
						R[x_1,\dots,x_n]/(f_1,\dots,f_m)\cong S
					\end{gather*}
				\end{itemize}
			\end{Def}
			\subsection{Ring Maps of Finite Type and of Finite Presentation}
			\subsection{Finite Ring Maps}
			\subsection{Homomorphisms Essentially of Finite Type}
				\begin{Def}
					Let $R\rightarrow S$ be a ring map:
					\begin{itemize}
						\item $R\rightarrow S$ is \textbf{essentially of finite type} if $S$ is the localization of an $R$-algebra of finite type. i.e.
						\begin{gather*}
							S=M^{-1}R[x_1,\dots,x_n]
						\end{gather*}
						\item $R\rightarrow S$ is \textbf{essentially of finite representation} if $S$ is the localization of an $R$-algebra of finite representation
					\end{itemize}
				\end{Def}
				
				\begin{lemma}
					The class of ring maps which are \textbf{essentially of finite type} is preserved under composition.
					
					Similarly for essentially finite representation.
				\end{lemma}
				\begin{proof}
					内容...
				\end{proof}
				\subsubsection{Under Base CHange}
					\begin{proposition}
						内容...
					\end{proposition}
					\begin{proof}
						内容...
					\end{proof}
				\subsubsection{Maps to Artinian Local Ring}
					\begin{lemma}
						内容...
					\end{lemma}
					\begin{proof}
						内容...
					\end{proof}
				\subsubsection{Essentially Finite Maps between Local Rings}
					\begin{lemma}
						If $\varphi:R\rightarrow S$ be essentially of finite type with $R$ and $S$ local(but, $\varphi$ may be not local) .
						
						Then, there exists an $n$ and a maximal ideal
						\begin{gather*}
							m\subseteq R[x_1,\dots,x_n]
						\end{gather*}
						
						lying over $m_R$ such that: $S$ is a localization of a quotient $R[x_1,\dots,x_n]_m$.
					\end{lemma}
					\begin{proof}
						内容...
					\end{proof}
	\chapter{Modules}
		\section{Internal Hom}
		\section{Tensor Product}
			
			\begin{Def}
			\end{Def}
			
			\begin{lemma}
				(Tensor Product commutes with Colimits) Let $(M_i,\mu_{ij})$ be a system over a preordered set $I$, then  
				
				\begin{gather*}
					colim(M_i\otimes N) \cong (Colim M_i)\otimes N
				\end{gather*}
			\end{lemma}
			
			\begin{proof}
			\end{proof}
			
			\begin{lemma}
				Let \begin{tikzcd}
					M_1\arrow[r,"f"]&M_2\arrow[r,"g"]&M_3\arrow[r]&0
				\end{tikzcd} be an exact sequence of $R$-modules and homomorphisms, and $N$ be an $R$-module, then the sequence
				
				\begin{center}
					\begin{tikzcd}
						M_1\otimes N\arrow[r,"f\otimes id"]&M_2\otimes N\arrow[r,"g\otimes id"]&M_3\otimes N\arrow[r]&0
					\end{tikzcd}
				\end{center}
				
				is exact. In other word, $-\otimes N$ is right exact.
			\end{lemma}
			
			\begin{proof}
			\end{proof}
			
			\subsection{Finiteness}
			
			\begin{lemma}
			\end{lemma}
			
			\begin{proof}
			\end{proof}
			
			\subsection{Tensor Algebra}
		
		\section{Limits, Colimits}
			\subsection{Directed System, Directed Limits}
			\begin{Def}
				Let $(I,\leq)$ be preordered set. Then, $(M_i,\mu_{ij})$ of $R$-modules over $I$ consists of a family of $R$-modules $\{M_i\}$ indexed by $I$, and a family of $R$-module maps $\{\mu_{ij}:M_i\rightarrow M_j\}_{i\leq j}$ such that for $i\leq j\leq k$ 
				\begin{gather*}
					\mu_{ii}=id_{M_i}\qquad \mu_{ik}=\mu_{jk}\circ \mu_{ij}
				\end{gather*}
				We say $(M_i,\mu_{ij})$ is a \textbf{directed system} if $I$ is a directed set.
			\end{Def}
			
			\begin{lemma}
				Let $(M_i,\mu_{ij})$ be a family of $R$-modules over the preordered set $I$. The colimit o the system $(M_i,\mu_{ij})$ is the quotient $R$-module
				\begin{gather*}
					(\oplus_i M_i)/Q\qquad Q=<\{\iota_i(x_i)-\iota_j(\mu_{ij}(x_i))\}_{i\leq j}>_R
				\end{gather*}
				
				where $\iota_i:M_i\rightarrow \oplus_{i\in I} M_i$ is a natural inclusion.
				
				We denote the colimit
				\begin{gather*}
					colimit(M)=colim_i M_i
				\end{gather*}
				
				Denote $\pi:\oplus_{i\in I} M_i\rightarrow M$ the projection map and $\phi_i=\pi\circ \iota_i:M_i\rightarrow M$.
			\end{lemma}
			\begin{proof}
				内容...
			\end{proof}
			
			\begin{lemma}
				Let $(M_i,\mu_{ij})$ be a system of $R$-modules over the preordered set $I$.
				
				Assume that $I$ is directed. The colimit of the system $(M_i.\mu_{ij})$ is canonically isomorphic to the module $M$ defined as follows:
				\begin{itemize}
					\item 
					\item 
					\item 
				\end{itemize}
			\end{lemma}
			\begin{proof}
				内容...
			\end{proof}
			
			\begin{lemma}
				Let $(M_i,\mu_{ij})$ be a directed system. Let $M=colim M_i$ with
				\begin{gather*}
					\mu_i:M_i\rightarrow M
				\end{gather*}
				Then, $\mu_i(x_i)=0,x_i\in M\Leftrightarrow $ There exists $j\geq i$ such that $\mu_{ij}(x_i)=0$
			\end{lemma}
			\begin{proof}
				内容...
			\end{proof}
			\subsubsection{Morphisms of Directed Systems}
			\begin{lemma}
				内容...
			\end{lemma}
			\begin{proof}
				内容...
			\end{proof}
			
			\begin{lemma}
				内容...
			\end{lemma}
			\begin{proof}
				内容...
			\end{proof}
			\subsection{Directed Limit in Index Category}
			
			\begin{proposition}
				内容...
			\end{proposition}
			\begin{proof}
				内容...
			\end{proof}
		\section{Characterizing Finite and Finitely Presented Modules}
			
		\section{Finite Case of Modules and Maps}
			\subsection{Finite Modules, and Finitely presented modules}
		
			\begin{Def}
				Let $R$ be a ring and $M$ is $R$-module:
				
				\begin{itemize}
					\item We say $M$ is a \textbf{finite} $R$-module, or finitely generated module if there exists $n\in \NN$ and $x_i\in M$, with exact sequence
					\begin{gather*}
						\oplus Rx_i\rightarrow M \rightarrow 0\quad\mbox{or}\quad R^{\oplus n}\rightarrow M\rightarrow 0
					\end{gather*}
					\item $M$ is \textbf{finitely represented} if there exists $m,n\in \NN$ and exact sequence
					\begin{gather*}
						R^{\oplus m}\rightarrow R^{\oplus n}\rightarrow M\rightarrow 0
					\end{gather*}
					which implies:
					\begin{itemize}
						\item $M$ is finitely generated;
						\item The relation between the generators are also finitely generated.
					\end{itemize}
					such choice of exact sequence is called a \textbf{presentation}.
 				\end{itemize}
			\end{Def}
			
			\begin{lemma}
				Let $R$ be a ring, and 
				
				\begin{gather*}
					\alpha:R^{\oplus n}\rightarrow M\qquad \beta:N\rightarrow M
				\end{gather*}
				
				be $R$-module maps.If $Im(\alpha)\subseteq Im(\beta)$, then there exists an $R$-module map $\gamma:R^{\oplus n}\rightarrow N$ such that $\alpha=\beta\circ\gamma$.
				
				\begin{center}
					\begin{tikzcd}
						N\arrow[r,"\beta"]&M\\
						R^{\oplus n}\arrow[u,dashrightarrow,"\gamma"]\arrow[ru,"\alpha"]
					\end{tikzcd}
				\end{center}
			\end{lemma}
			
			\begin{proof}
				Let $\alpha(e_i)=\beta(n_i)$. Then, consider the mapping
				
				\begin{gather*}
					\gamma:R^{\oplus n}\rightarrow N\qquad (a_1,\dots,a_n)\mapsto \sum a_i n_i
				\end{gather*}
			\end{proof}
			
			\begin{lemma}
				Let $R$ be a ring, and $0\rightarrow M_1\rightarrow M_2\rightarrow M_3\rightarrow 0$ be a short exact sequence of $R$-modules:
				
				\begin{itemize}
					\item If $M_1$ and $M_3$ are finite $R$-modules, then $M_2$ is a finite $R$-module;
					\item If $M_1$ and $M_3$ are finitely presented $R$-modules, then $M_2$ is finitely presented;
					\item If $M_2$ is a finite $R$-module, then $M_3$ is a finite $R$-module;
					\item If $M_2$ is finitely presented $R$-module and $M_1$ is a finite $R$-module, then $M_3$ is a finitely represented $R$-module;
					\item If $M_3$ is finitely presented $R$-module and $M_2$ is a finite $R$-module,then $M_1$ is a finite $R$-module. 
				\end{itemize}
			\end{lemma}
			
			\begin{proof}
				\begin{itemize}
					\item 
					\item 
					\item 
					\item 
					\item 
				\end{itemize}
			\end{proof}
			
			\begin{lemma}
				Let $R$ be a ring and let $M$ be a finite $R$-module. There exists a filtration by finite $R$-module 
				\begin{gather*}
					0\subseteq M_0 \subseteq M_1\subseteq \dots\subseteq M_n=M
				\end{gather*}
				
				such that each quotient $M_i/M_{i-1}$ is isomorphic to $R/I_i$ for some ideal $I_i$ of $R$.
			\end{lemma}
			
			\begin{proof}
				内容...
			\end{proof}
			
			\begin{lemma}
				Let $R\rightarrow S$ be a map, and $M$ be an $S$-module. If $M$ is finite as an $R$-module, then $M$ is finite as an $S$-module.	
			\end{lemma}
			
			\begin{proof}
				内容...
			\end{proof}
		\section{Base Change}
	\chapter{Associative Algebra}
	
	\chapter{Polynomial Algebra}
		\section{Discriminant }
	\chapter{Matrix Algebra}
\end{document}